% AA-42 | Reprezentacje zespolone grup skonczonych: Schur, charaktery, ortogonalnosc.
\chapter{Lemat Schura, charaktery i relacje ortogonalności}
\markboth{Lemat Schura i charaktery}{Lemat Schura i charaktery}
\label{aa:ch:schur-characters}

\section{Dlaczego same rozkłady reprezentacji nie wystarczą?}
Twierdzenie Maschkego z Rozdziału~\ref{aa:ch:representations}
zapewnia rozkład każdej skończenie wymiarowej reprezentacji
grupy skończonej na składniki nieprzywiedlne, ale samo nie daje
sposobu obliczenia ich krotności ani rozpoznawania izomorfizmów.
Narzędziami służącymi do tego celu są operatory przeplatające,
lemat Schura i charaktery. Przy badaniu relacji ortogonalności
będziemy pracować nad ciałem \(\mathbb C\); warunek ten
jest istotny i nie może zostać bez komentarza przeniesiony
na dowolne ciało.

\begin{definicja}[Operator przeplatający i równoważność reprezentacji]
\label{aa:def:intertwiner}
Niech \(G\) będzie grupą, \(K\) ciałem, a
\(\rho:G\to\operatorname{GL}_K(V)\) oraz
\(\sigma:G\to\operatorname{GL}_K(W)\) reprezentacjami
na przestrzeniach wektorowych \(V,W\) nad \(K\).
Operatorem przeplatającym te reprezentacje nazywamy
odwzorowanie \(K\)-liniowe \(T:V\to W\), dla którego
\(T\rho(g)=\sigma(g)T\) dla każdego \(g\in G\).
Przestrzeń wszystkich takich operatorów oznaczamy przez
\(\operatorname{Hom}_G(V,W)\).
Reprezentacje są równoważne, gdy istnieje
odwracalny operator z \(\operatorname{Hom}_G(V,W)\).
\end{definicja}

\begin{lemat}[Pierwszy lemat Schura]\label{aa:lem:schur-first}
Niech \(G\) będzie grupą i \(K\) ciałem.
Jeżeli \(V\) i \(W\) są nieprzywiedlnymi, skończenie
wymiarowymi reprezentacjami \(G\) nad \(K\), to każdy
niezerowy \(T\in\operatorname{Hom}_G(V,W)\) jest izomorfizmem.
W szczególności, jeżeli \(V\) i \(W\) nie są równoważne,
to \(\operatorname{Hom}_G(V,W)=\{0\}\).
\end{lemat}
\begin{proof}
Jeżeli \(v\in\ker T\), to
\(T\rho(g)v=\sigma(g)Tv=0\), więc jądro jest
podprzestrzenią \(G\)-niezmienniczą przestrzeni \(V\).
Analogicznie \(\sigma(g)T(V)=T\rho(g)(V)\subseteq T(V)\),
zatem obraz jest \(G\)-niezmienniczy w \(W\).
Nieprzywiedlność daje \(\ker T=0\) albo \(\ker T=V\),
oraz \(T(V)=0\) albo \(T(V)=W\).
Dla \(T\ne0\) zachodzą odpowiednio \(\ker T=0\)
i \(T(V)=W\), co oznacza, że \(T\) jest izomorfizmem.
\end{proof}

\begin{twierdzenie}[Lemat Schura nad \(\mathbb C\)]
\label{aa:tw:schur-complex}
Niech \(G\) będzie grupą, a \(V\ne0\) skończenie
wymiarową, nieprzywiedlną reprezentacją \(G\)
nad \(\mathbb C\). Wówczas każdy operator
\(T\in\operatorname{Hom}_G(V,V)\) jest postaci
\(T=\lambda I_V\) dla pewnej liczby
\(\lambda\in\mathbb C\).
\end{twierdzenie}
\begin{proof}
Wielomian charakterystyczny operatora \(T\)
ma dodatni stopień, więc na mocy zasadniczego
twierdzenia algebry posiada pierwiastek
\(\lambda\in\mathbb C\).
Operator \(T-\lambda I_V\) przeplata działanie \(G\)
i nie jest odwracalny, gdyż ma niezerowe jądro.
Z pierwszego lematu Schura wynika, że gdyby był
niezerowy, musiałby być izomorfizmem.
Dlatego \(T-\lambda I_V=0\).
\end{proof}

\begin{przyklad}[Dlaczego ciało \(\mathbb R\) nie wystarcza]
\label{aa:prz:schur-real}
Niech \(G=C_3=\langle g\rangle\) i niech generator działa
na \(\mathbb R^2\) przez obrót o \(2\pi/3\):
\[
\rho(g)=
\begin{pmatrix}
-\frac12&-\frac{\sqrt3}{2}\\
\frac{\sqrt3}{2}&-\frac12
\end{pmatrix}.
\]
Macierz ma wielomian charakterystyczny \(T^2+T+1\),
bez pierwiastków rzeczywistych. Nie istnieje więc
prosta niezmiennicza, a reprezentacja jest nieprzywiedlna
nad \(\mathbb R\). Operator
\(J=\left(\begin{smallmatrix}0&-1\\1&0\end{smallmatrix}\right)\)
komutuje z \(\rho(g)\), lecz nie jest wielokrotnością
\(I\). Zatem końcowa teza lematu Schura nad \(\mathbb C\)
nie zachodzi w tej postaci nad \(\mathbb R\).
\end{przyklad}

\section{Uśrednianie iloczynu hermitowskiego}
\begin{stwierdzenie}[Unitarność reprezentacji grupy skończonej]
\label{aa:prop:unitarization}
Niech \(G\) będzie grupą skończoną, a
\(\rho:G\to\operatorname{GL}_{\mathbb C}(V)\)
reprezentacją na skończenie wymiarowej przestrzeni
zespolonej \(V\). Istnieje dodatnio określony
iloczyn hermitowski na \(V\), względem którego
wszystkie \(\rho(g)\) są unitarne.
\end{stwierdzenie}
\begin{proof}
Wybierzmy dowolny dodatnio określony iloczyn hermitowski
\(\langle\cdot,\cdot\rangle_0\) na \(V\) i połóżmy
\[
\langle v,w\rangle_G
=\frac1{|G|}\sum_{g\in G}
\langle\rho(g)v,\rho(g)w\rangle_0.
\]
Skończona suma zachowuje własności
półtoraliniowości oraz symetrii hermitowskiej.
Dla \(v\ne0\) każdy składnik
\(\langle\rho(g)v,\rho(g)v\rangle_0\) jest dodatni,
bo \(\rho(g)\) jest odwracalne. Iloczyn jest zatem
dodatnio określony. Dla \(h\in G\) podstawienie
\(k=gh\) daje
\(\langle\rho(h)v,\rho(h)w\rangle_G
=\langle v,w\rangle_G\).
Z definicji oznacza to unitarność każdego \(\rho(h)\).
\end{proof}

\begin{definicja}[Charakter reprezentacji]
\label{aa:def:character}
Niech \(G\) będzie grupą skończoną, a
\(\rho:G\to\operatorname{GL}_{\mathbb C}(V)\)
reprezentacją na skończenie wymiarowej przestrzeni
zespolonej \(V\). Jej charakterem nazywamy funkcję
\(\chi_\rho:G\to\mathbb C\) określoną wzorem
\(\chi_\rho(g)=\operatorname{tr}(\rho(g))\).
Charakter reprezentacji nieprzywiedlnej nazywamy
charakterem nieprzywiedlnym.
Funkcję \(f:G\to\mathbb C\) nazywamy funkcją klasową,
gdy \(f(hgh^{-1})=f(g)\) dla wszystkich \(g,h\in G\).
\end{definicja}

\begin{stwierdzenie}[Podstawowe własności charakterów]
\label{aa:prop:character-basics}
Niech \(G\) będzie grupą skończoną, a
\(\rho:G\to\operatorname{GL}_{\mathbb C}(V)\)
skończenie wymiarową reprezentacją zespoloną
o charakterze \(\chi_\rho\). Wtedy
\(\chi_\rho(e)=\dim_{\mathbb C}V\),
\(\chi_\rho(hgh^{-1})=\chi_\rho(g)\) oraz
\(\chi_\rho(g^{-1})=\overline{\chi_\rho(g)}\).
Ponadto dla reprezentacji \(\rho_1,\rho_2\)
zachodzi \(\chi_{\rho_1\oplus\rho_2}
=\chi_{\rho_1}+\chi_{\rho_2}\).
\end{stwierdzenie}
\begin{proof}
Pierwsza równość wynika z \(\rho(e)=I_V\).
Ponieważ
\(\rho(hgh^{-1})=\rho(h)\rho(g)\rho(h)^{-1}\),
niezmienniczość śladu przy podobieństwie daje drugą.
Możemy przyjąć iloczyn hermitowski ze
Stwierdzenia~\ref{aa:prop:unitarization}.
Wtedy \(\rho(g^{-1})=\rho(g)^*\), a ślad operatora
sprzężonego hermitowsko jest sprzężeniem śladu
wyjściowego. Ostatnia własność wynika z postaci
blokowo diagonalnej macierzy sumy prostej.
\end{proof}

\section{Relacje ortogonalności}
\begin{definicja}[Iloczyn skalarny funkcji na grupie]
\label{aa:def:character-inner-product}
Niech \(G\) będzie grupą skończoną.
Dla funkcji \(f,h:G\to\mathbb C\) określamy
\[
\langle f,h\rangle_G
=\frac1{|G|}\sum_{g\in G}f(g)\overline{h(g)}.
\]
Jest to dodatnio określony iloczyn hermitowski
na zespolonej przestrzeni wszystkich funkcji \(G\to\mathbb C\).
Ogranicza się on do przestrzeni funkcji klasowych.
\end{definicja}

\begin{twierdzenie}[Ortogonalność współczynników macierzowych]
\label{aa:tw:matrix-orthogonality}
Niech \(G\) będzie grupą skończoną oraz
\(\rho_i:G\to U(V_i)\), \(1\le i\le r\),
skończenie wymiarowymi nieprzywiedlnymi reprezentacjami
zespolonymi, parami nierównoważnymi.
Niech \(d_i=\dim_{\mathbb C}V_i\) i wybierzmy
bazę ortonormalną każdej \(V_i\). Dla współczynników
macierzy \(\rho_i(g)\) zachodzi
\[
\frac1{|G|}\sum_{g\in G}
(\rho_i(g))_{ab}\overline{(\rho_j(g))_{cd}}
=\begin{cases}
\dfrac{\delta_{ac}\delta_{bd}}{d_i},&i=j,\\[3pt]
0,&i\ne j,
\end{cases}
\]
dla \(1\le a,b\le d_i\) oraz \(1\le c,d\le d_j\).
\end{twierdzenie}
\begin{proof}
Dla liniowego \(A:V_j\to V_i\) rozważmy operator
\[
\mathcal P_{ij}(A)=\frac1{|G|}\sum_{g\in G}
\rho_i(g)A\rho_j(g)^{-1}.
\]
Podstawienie \(k=hg\) pokazuje, że
\(\rho_i(h)\mathcal P_{ij}(A)
=\mathcal P_{ij}(A)\rho_j(h)\) dla każdego \(h\in G\).
Gdy \(i\ne j\), pierwszy lemat Schura daje
\(\mathcal P_{ij}(A)=0\). Gdy \(i=j\), zespolony
lemat Schura zapewnia \(\mathcal P_{ii}(A)=\lambda I\).
Z niezmienniczości śladu przy sprzężeniu otrzymujemy
\(d_i\lambda=\operatorname{tr}A\), czyli
\(\mathcal P_{ii}(A)=(\operatorname{tr}A/d_i)I\).

Niech teraz \(A\) wysyła \(d\)-ty wektor bazy \(V_j\)
na \(b\)-ty wektor bazy \(V_i\), a pozostałe
wektory bazowe na zero. Wówczas współczynnik
\((a,c)\) operatora \(\mathcal P_{ij}(A)\) wynosi
\[
\frac1{|G|}\sum_{g\in G}
(\rho_i(g))_{ab}(\rho_j(g)^{-1})_{dc}
=\frac1{|G|}\sum_{g\in G}
(\rho_i(g))_{ab}\overline{(\rho_j(g))_{cd}},
\]
gdzie ostatnia równość korzysta z unitarności.
Gdy \(i=j\), mamy
\(\operatorname{tr}A=\delta_{bd}\), a współczynnik
\((a,c)\) operatora \(I\) jest równy \(\delta_{ac}\).
Dostajemy obie postacie żądanej relacji.
\end{proof}

\begin{wniosek}[Ortogonalność charakterów nieprzywiedlnych]
\label{aa:cor:character-orthogonality}
Niech \(G\) będzie grupą skończoną, a
\(\chi_1,\ldots,\chi_r\) charakterami parami
nierównoważnych, skończenie wymiarowych
nieprzywiedlnych reprezentacji zespolonych \(G\).
Wtedy \(\langle\chi_i,\chi_j\rangle_G=\delta_{ij}\).
\end{wniosek}
\begin{proof}
Dzięki unitarności możemy zastosować
Twierdzenie~\ref{aa:tw:matrix-orthogonality}.
Wyrażamy oba ślady jako sumy elementów diagonalnych:
\[
\langle\chi_i,\chi_j\rangle_G
=\sum_{a=1}^{d_i}\sum_{c=1}^{d_j}
\frac1{|G|}\sum_{g\in G}
(\rho_i(g))_{aa}\overline{(\rho_j(g))_{cc}}.
\]
Dla \(i\ne j\) każdy składnik jest zerowy.
Dla \(i=j\) wewnętrzna suma wynosi
\(\delta_{ac}/d_i\); suma po \(a,c\) jest równa \(1\).
\end{proof}

\section{Krotności składników i reprezentacja regularna}
\begin{stwierdzenie}[Występowanie nieprzywiedlnych składników]
\label{aa:prop:regular-contains-all}
Niech \(G\) będzie grupą skończoną.
Każda skończenie wymiarowa nieprzywiedlna reprezentacja
zespolona \(V\) grupy \(G\) jest równoważna składnikowi
reprezentacji regularnej na \(\mathbb C[G]\).
W szczególności istnieje tylko skończenie wiele klas
równoważności takich reprezentacji.
\end{stwierdzenie}
\begin{proof}
Wybierzmy \(0\ne v\in V\) i określmy
\(G\)-równoważne odwzorowanie
\(\Phi:\mathbb C[G]\to V\) przez
\(\Phi([g])=\rho(g)v\), przedłużając liniowo.
Jego obraz jest niezerową podprzestrzenią niezmienniczą,
a więc dzięki nieprzywiedlności jest równy \(V\).
Twierdzenie Maschkego zapewnia niezmiennicze dopełnienie
\(\ker\Phi\) w \(\mathbb C[G]\); ograniczenie \(\Phi\)
do tego dopełnienia daje izomorfizm z \(V\).
Ponadto rozkład skończenie wymiarowej reprezentacji
regularnej na nieprzywiedlne składniki ma skończenie
wiele wyrazów. Każda klasa nieprzywiedlna pojawia
się w tym rozkładzie: w istocie izomorficzny z \(V\)
podmoduł, który przed chwilą skonstruowaliśmy,
ma charakter zawarty w rozkładzie regularnym.
Ściślej, jeśli \(\mathbb C[G]=\bigoplus_j S_j\),
to rzut na \(V\) złożony z włączeniami \(S_j\)
musi być niezerowy dla pewnego \(j\). Pierwszy
lemat Schura daje wtedy \(S_j\simeq V\).
Liczba różnych klas jest więc skończona.
\end{proof}

\begin{twierdzenie}[Krotności z iloczynów skalarnych]
\label{aa:tw:multiplicities}
Niech \(G\) będzie grupą skończoną.
Niech \(V_1,\ldots,V_r\) będą przedstawicielami
wszystkich klas równoważności nieprzywiedlnych
reprezentacji zespolonych \(G\), z charakterami
\(\chi_1,\ldots,\chi_r\). Jeżeli \(V\) jest dowolną
skończenie wymiarową reprezentacją zespoloną
o charakterze \(\chi_V\), to istnieją jednoznaczne
liczby całkowite \(m_i\ge0\) takie, że
\[
V\simeq\bigoplus_{i=1}^r V_i^{\oplus m_i},
\qquad m_i=\langle\chi_V,\chi_i\rangle_G.
\]
W konsekwencji dwie skończenie wymiarowe reprezentacje
zespolone grupy \(G\) są równoważne wtedy i tylko
wtedy, gdy mają te same charaktery.
\end{twierdzenie}
\begin{proof}
Istnienie rozkładu wynika z twierdzenia Maschkego,
a kompletność listy \(V_i\) ze
Stwierdzenia~\ref{aa:prop:regular-contains-all}.
Addytywność śladu daje
\(\chi_V=\sum_i m_i\chi_i\).
Iloczyn skalarny z \(\chi_j\) i ortogonalność
charakterów prowadzą do
\(\langle\chi_V,\chi_j\rangle_G=m_j\).
Stąd krotności są wyznaczone jednoznacznie.
Jeśli charaktery dwóch reprezentacji są równe,
ich krotności wszystkich składników są jednakowe,
więc obie reprezentacje są równoważne.
Odwrotność wynika z niezmienniczości śladu
przy podobieństwie.
\end{proof}

\begin{twierdzenie}[Charakter regularny i suma kwadratów wymiarów]
\label{aa:tw:regular-character}
Niech \(G\) będzie grupą skończoną i niech
\(\rho_{\mathrm{reg}}\) będzie lewą reprezentacją
regularną na \(\mathbb C[G]\). Jej charakter spełnia
\[
\chi_{\mathrm{reg}}(g)=
\begin{cases}|G|,&g=e,\\0,&g\ne e.\end{cases}
\]
Jeżeli \(d_i=\dim_{\mathbb C}V_i=\chi_i(e)\)
dla pełnej rodziny \(V_1,\ldots,V_r\) jak wyżej,
to \(\mathbb C[G]\simeq\bigoplus_{i=1}^r
V_i^{\oplus d_i}\) oraz
\(\sum_{i=1}^r d_i^2=|G|\).
\end{twierdzenie}
\begin{proof}
W bazie \(([h])_{h\in G}\) działanie elementu \(g\)
permutuje bazę według \([h]\mapsto[gh]\).
Jeżeli \(g\ne e\), to żaden wektor bazowy nie
jest stały, bo \(gh=h\) wymusza \(g=e\).
Ślad macierzy permutacyjnej jest liczbą
nieruchomych wektorów bazowych, zatem wynosi \(0\);
dla \(g=e\) wynosi \(|G|\).
Z Twierdzenia~\ref{aa:tw:multiplicities} krotność
\(V_i\) w reprezentacji regularnej to
\[
\langle\chi_{\mathrm{reg}},\chi_i\rangle_G
=\frac1{|G|}|G|\overline{\chi_i(e)}=d_i.
\]
Porównując wymiary obu stron otrzymanego rozkładu,
dostajemy \(|G|=\sum_i d_i^2\).
\end{proof}

\begin{stwierdzenie}[Jawny rzut na składnik izotypiczny]
\label{aa:prop:isotypic-projector}
Niech \(G\) będzie grupą skończoną, a
\(V\) skończenie wymiarową reprezentacją zespoloną
o działaniu \(\rho\). Dla każdego charakteru
nieprzywiedlnego \(\chi_i\), którego wymiar
wynosi \(d_i=\chi_i(e)\), określmy
\[
P_i=\frac{d_i}{|G|}\sum_{g\in G}
\overline{\chi_i(g)}\,\rho(g).
\]
Wtedy \(P_i\) jest \(G\)-równoważnym rzutem
na sumę wszystkich składników \(V\) równoważnych
\(V_i\) w dowolnym rozkładzie na składniki
nieprzywiedlne. Ponadto \(P_iP_j=0\) dla \(i\ne j\)
oraz \(\sum_{i=1}^rP_i=I_V\), gdy \(\chi_1,\ldots,\chi_r\)
stanowią pełną rodzinę.
\end{stwierdzenie}
\begin{proof}
Charakter \(\chi_i\) jest funkcją klasową, zatem
sprzężenie sumy definiującej \(P_i\) przez
\(\rho(h)\) jedynie permutuje jej składniki.
Wobec tego \(P_i\) przeplata działanie \(G\).
Z twierdzenia Maschkego weźmy rozkład
\(V=\bigoplus_\alpha W_\alpha\) na nieprzywiedlne
podreprezentacje. Każda z nich jest zachowana przez
\(P_i\), ponieważ sumujemy operatory \(\rho(g)\),
które osobno zachowują \(W_\alpha\).
Gdy \(W_\alpha\simeq V_j\), zespolony lemat Schura
daje \(P_i|_{W_\alpha}=\lambda I\). Obliczając ślad,
otrzymujemy
\[
d_j\lambda
=\frac{d_i}{|G|}\sum_{g\in G}
\overline{\chi_i(g)}\chi_j(g)
=d_i\delta_{ij}.
\]
Jeżeli \(j=i\), to \(\lambda=1\); w przeciwnym
przypadku \(\lambda=0\). To ustala działanie
operatorów \(P_i\) na każdym składniku rozkładu
i dowodzi wszystkich wskazanych własności.
\end{proof}

\section{Funkcje klasowe i ortogonalność kolumn}
\begin{twierdzenie}[Charaktery jako baza funkcji klasowych]
\label{aa:tw:characters-basis}
Niech \(G\) będzie grupą skończoną.
Wszystkie charaktery nieprzywiedlne zespolone
\(\chi_1,\ldots,\chi_r\) tworzą bazę ortonormalną
przestrzeni funkcji klasowych \(G\to\mathbb C\)
względem iloczynu skalarnego
z Definicji~\ref{aa:def:character-inner-product}.
W szczególności liczba \(r\) klas równoważności
nieprzywiedlnych reprezentacji zespolonych \(G\)
jest równa liczbie klas sprzężoności w \(G\).
\end{twierdzenie}
\begin{proof}
Charaktery są funkcjami klasowymi i są ortonormalne
na mocy Wniosku~\ref{aa:cor:character-orthogonality}.
Aby wykazać zupełność, załóżmy, że funkcja klasowa
\(f:G\to\mathbb C\) jest prostopadła do każdego
\(\chi_i\). Dla dowolnej reprezentacji \(\rho\)
grupy \(G\) rozważmy operator
\[
T_f^\rho=\sum_{g\in G}f(g)\rho(g^{-1}).
\]
Ponieważ \(f\) jest stała na klasach sprzężoności,
operator ten przeplata działanie \(G\).
Na nieprzywiedlnej reprezentacji \(V_i\)
lemat Schura daje \(T_f^{\rho_i}=\lambda_iI\).
Biorąc ślad i korzystając z równości
\(\chi_i(g^{-1})=\overline{\chi_i(g)}\), mamy
\[
d_i\lambda_i=\sum_{g\in G}
f(g)\overline{\chi_i(g)}
=|G|\langle f,\chi_i\rangle_G=0.
\]
Zatem \(T_f^{\rho_i}=0\) dla wszystkich \(i\).
Reprezentacja regularna jest sumą takich składników,
więc \(T_f^{\mathrm{reg}}=0\). Jednak
\(T_f^{\mathrm{reg}}([e])
=\sum_{g\in G}f(g)[g^{-1}]\).
Liniowa niezależność formalnej bazy \(([g])\)
wymusza \(f(g)=0\) dla każdego \(g\in G\).
W przestrzeni skończenie wymiarowej brak
niezerowego wektora prostopadłego do rozpiętości
ortonormalnych charakterów oznacza, że rozpinają
całą przestrzeń funkcji klasowych.
Jej wymiar jest liczbą klas sprzężoności, ponieważ
funkcja klasowa jest jednoznacznie wyznaczona przez
jedną wartość na każdej klasie. Stąd \(r\) jest równy
tej liczbie.
\end{proof}

\begin{wniosek}[Ortogonalność kolumn tablicy charakterów]
\label{aa:cor:column-orthogonality}
Niech \(G\) będzie grupą skończoną,
\(C_1,\ldots,C_r\) jej wszystkimi klasami sprzężoności
i \(g_k\in C_k\) ustalonymi reprezentantami.
Jeżeli \(\chi_1,\ldots,\chi_r\) są wszystkimi
charakterami nieprzywiedlnymi zespolonymi, to
\[
\sum_{i=1}^r\chi_i(g_k)\overline{\chi_i(g_\ell)}
=\begin{cases}
|G|/|C_k|,&k=\ell,\\
0,&k\ne\ell.
\end{cases}
\]
\end{wniosek}
\begin{proof}
W przestrzeni funkcji klasowych iloczyn skalarny ma
postać
\(\langle f,h\rangle_G
=\sum_{k=1}^r (|C_k|/|G|)f(g_k)\overline{h(g_k)}\).
Z Twierdzenia~\ref{aa:tw:characters-basis}
macierz kwadratowa
\(A_{ik}=\sqrt{|C_k|/|G|}\,\chi_i(g_k)\)
ma ortonormalne wiersze. Jest więc unitarna,
a jej kolumny także są ortonormalne.
Równości \((A^*A)_{k\ell}=\delta_{k\ell}\)
dają tezę po usunięciu dodatnich czynników
\(\sqrt{|C_k||C_\ell|}/|G|\).
\end{proof}

\begin{wniosek}[Rozwinięcie funkcji klasowej]
\label{aa:cor:class-fourier}
Niech \(G\) będzie grupą skończoną, a
\(f:G\to\mathbb C\) funkcją klasową.
Dla pełnej rodziny charakterów
nieprzywiedlnych \(\chi_1,\ldots,\chi_r\) zachodzi
\[
f=\sum_{i=1}^r\langle f,\chi_i\rangle_G\,\chi_i.
\]
Współczynniki są wyznaczone jednoznacznie.
\end{wniosek}
\begin{proof}
Twierdzenie~\ref{aa:tw:characters-basis} zapewnia
rozwinięcie w bazie ortonormalnej. Biorąc iloczyn
skalarny obu stron z \(\chi_j\), dostajemy
jednoznacznie współczynnik przy \(\chi_j\).
\end{proof}

\section{Dwie tablice charakterów}
\begin{przyklad}[Wszystkie charaktery grupy \(C_3\)]
\label{aa:prz:character-c3}
Niech \(C_3=\langle g:g^3=e\rangle\) oraz
\(\zeta=\exp(2\pi i/3)\). Trzy reprezentacje
jednowymiarowe określone wzorami
\(\rho_j(g)=\zeta^j\), \(j=0,1,2\), mają charaktery
\(\chi_j(g^k)=\zeta^{jk}\).
Są nierównoważne, gdyż mają różne wartości na \(g\).
Równość
\(\sum_{k=0}^2\zeta^{(j-\ell)k}
=3\delta_{j\ell}\) dowodzi bezpośrednio ich
ortogonalności. Wzór na sumę kwadratów wymiarów
daje \(1^2+1^2+1^2=3\), a zatem nie występują
żadne inne klasy reprezentacji nieprzywiedlnych.
Tablica charakterów ma postać
\[
\begin{array}{c|ccc}
 & e & g & g^2\\ \hline
\chi_0&1&1&1\\
\chi_1&1&\zeta&\zeta^2\\
\chi_2&1&\zeta^2&\zeta
\end{array}.
\]
Ponieważ \(C_3\) jest abelowa, każda klasa
sprzężoności jest jednoelementowa. Ortogonalność
kolumn wynika stąd również przez bezpośrednie
obliczenie iloczynów kolumn tej tablicy.
\end{przyklad}

\begin{przyklad}[Pełna tablica charakterów \(S_3\)]
\label{aa:prz:character-s3}
Grupa \(S_3\) ma trzy klasy sprzężoności:
jednostkę \(e\), trzy transpozycje \(t\) i dwa
trójcykle \(c\). Reprezentacja trywialna ma charakter
\((1,1,1)\), a reprezentacja znaku \((1,-1,1)\).
Dla permutacyjnej reprezentacji \(\mathbb C^3\)
ślad permutacji jest liczbą nieruchomych wektorów
bazy. Jej charakter przyjmuje więc wartości
\((3,1,0)\). Odejmując charakter składnika
trywialnego, otrzymujemy charakter
\(\chi_{\mathrm{std}}=(2,0,-1)\) reprezentacji
na płaszczyźnie \(x_1+x_2+x_3=0\).

W tej sytuacji dla funkcji klasowych
\(\langle f,h\rangle_{S_3}
=(f(e)\overline{h(e)}
+3f(t)\overline{h(t)}
+2f(c)\overline{h(c)})/6\).
Stąd
\(\langle\chi_{\mathrm{std}},
\chi_{\mathrm{std}}\rangle_{S_3}
=(4+2)/6=1\),
więc reprezentacja standardowa jest nieprzywiedlna:
gdyby była sumą co najmniej dwóch nieprzywiedlnych,
jej iloczyn skalarny z samą sobą byłby sumą
kwadratów dodatnich całkowitych krotności,
a więc wynosiłby co najmniej \(2\).
Iloczyny mieszane trzech charakterów są zerowe,
co można sprawdzić na przykładzie
\((2+2(-1))/6=0\) dla charakteru trywialnego
i standardowego.

Pełna tablica jest następująca:
\[
\begin{array}{c|ccc}
& [e]\ (1)&[t]\ (3)&[c]\ (2)\\ \hline
\chi_{\mathrm{triv}}&1&1&1\\
\chi_{\mathrm{sgn}}&1&-1&1\\
\chi_{\mathrm{std}}&2&0&-1
\end{array}.
\]
Liczby w nawiasach to rozmiary klas. Ponieważ
\(1^2+1^2+2^2=6=|S_3|\), tablica obejmuje
wszystkie reprezentacje nieprzywiedlne.
W pierwszej kolumnie suma kwadratów to \(6\),
a w kolumnach transpozycji i trójcykli
wynosi odpowiednio \(2=6/3\) i \(3=6/2\),
zgodnie z ortogonalnością kolumn.
\end{przyklad}

\begin{uwaga}[Granice stosowalności]
\label{aa:rem:characters-limits}
Wnioski o ortonormalności charakterów,
rozpoznawaniu reprezentacji po ich charakterach
i sumie kwadratów wymiarów zostały udowodnione
dla reprezentacji zespolonych grup skończonych.
Nie wolno ich bezpośrednio stosować w
charakterystyce dzielącej rząd grupy:
wtedy twierdzenie Maschkego może zawieść,
a różne nierównoważne reprezentacje mogą mieć
ten sam charakter śladowy. Przykładowo nad
\(\mathbb F_2\) trywialna reprezentacja \(C_2\)
w wymiarze \(2\) i nietrywialne działanie
\(s\mapsto\left(\begin{smallmatrix}1&1\\0&1\end{smallmatrix}\right)\)
mają ślad \(0\) dla obu elementów \(C_2\),
chociaż reprezentacje nie są równoważne:
pierwsza ma działanie trywialne, druga nie.
\end{uwaga}
